\chapter{Overall Conclusions} \label{sec:cp5}
\newpage

KEEP CALM and CARRY ON.\\It's almost over!

On the operational side, the integration of the DA system into real-time forecasts represents the ultimate goal. In this context, improvements in computational efficiency and infrastructure will be critical, as they could enable the adoption of more sophisticated DA methods for both analysis and reanalysis applications. Furthermore, although SMS-Coastal provides a robust foundation for operational forecasting, its full integration with the SOMA assimilation system, as well as its validation across additional MOHID-based applications, remain important steps for future development. From this perspective, the present work should be understood not as a final solution, but as a foundational contribution that establishes the basis for the systematic integration of DA methodologies in the Algarve region. By developing and validating the SOMA OSSE System, demonstrating its application to observation strategy design, and providing the operational infrastructure required to support these processes, this thesis opens new pathways for advancing both ocean observation and numerical modelling in the region. As such, it lays the groundwork for future developments that will enhance forecast capabilities, support more informed decision-making, and contribute to the evolution of operational oceanography in coastal environments.
